% =====================================================================
% Seccion: MVP de demostracion (Hito 2)
% Fragmento SIN preambulo: % =====================================================================
% Seccion: MVP de demostracion (Hito 2)
% Fragmento SIN preambulo: % =====================================================================
% Seccion: MVP de demostracion (Hito 2)
% Fragmento SIN preambulo: % =====================================================================
% Seccion: MVP de demostracion (Hito 2)
% Fragmento SIN preambulo: \input{secciones/mvp}
% Capturas tomadas en el commit 58b1a9e, ultimo con el MVP disponible.
% =====================================================================
\section{MVP de demostración}
\label{sec:mvp}

Para el Hito~2 se construyó un MVP local, que se clona y se ejecuta con un solo comando: un backend FastAPI con los cuatro modelos y una interfaz estática en \arch{/mvp}. Su objetivo era demostrar la simulación de resultados y la trazabilidad de una predicción. Las capturas de esta sección se tomaron en el commit \arch{58b1a9e}. Tras presentar el hito, el MVP se retiró del proyecto (Sección~\ref{sec:arreglos}); su código sigue disponible en el historial de Git.

El MVP tenía cinco pestañas: caso a caso, subir imagen, simulación, modelos y métricas, y estudio de resonancia. Esta última la agregó Juan Manuel Tengonó y fue la base de la aplicación final.

\paragraph{Caso a caso y trazabilidad.} Se elige uno de los 45 casos de prueba (15 por clase) y se ejecutan los cuatro modelos. La Figura~\ref{fig:mvpcaso} muestra un error real: en \arch{caso_04}, con etiqueta Degeneración, los cuatro modelos predicen Protrusión. La Figura~\ref{fig:mvptraza} muestra el recorrido de esa imagen con valores de la ejecución: la ResNet-50 se calcula una sola vez y sus características pasan a los modelos 2, 3 y 4.

\begin{figure}[htbp]
\centering
\includegraphics[width=\linewidth]{mvp_caso.jpg}
\caption{Pestaña Caso a caso del MVP: resultado de los cuatro modelos sobre un caso de prueba.}
\label{fig:mvpcaso}
\end{figure}

\begin{figure}[htbp]
\centering
\includegraphics[width=\linewidth]{mvp_trazabilidad.jpg}
\caption{Trazabilidad de una imagen por los cuatro modelos, con tiempos y valores intermedios.}
\label{fig:mvptraza}
\end{figure}

La pestaña \emph{Simulación} se describe aparte, en la Sección~\ref{sec:simulacion}.

\paragraph{Modelos y métricas.} La pestaña reunía las métricas oficiales, las matrices de confusión, las curvas ROC y el SHA-256 de cada archivo de pesos (Figuras~\ref{fig:mvpmetricas} y~\ref{fig:mvproc}). Esta información ya no está en la aplicación: se conserva en la documentación y en la Sección~\ref{sec:resultados}.

\begin{figure}[htbp]
\centering
\includegraphics[width=\linewidth]{mvp_metricas.jpg}
\caption{Pestaña Modelos y métricas del MVP.}
\label{fig:mvpmetricas}
\end{figure}

\begin{figure}[htbp]
\centering
\includegraphics[width=\linewidth]{mvp_roc.jpg}
\caption{Curvas ROC y barras de exactitud y F1 macro en el MVP.}
\label{fig:mvproc}
\end{figure}

% Capturas tomadas en el commit 58b1a9e, ultimo con el MVP disponible.
% =====================================================================
\section{MVP de demostración}
\label{sec:mvp}

Para el Hito~2 se construyó un MVP local, que se clona y se ejecuta con un solo comando: un backend FastAPI con los cuatro modelos y una interfaz estática en \arch{/mvp}. Su objetivo era demostrar la simulación de resultados y la trazabilidad de una predicción. Las capturas de esta sección se tomaron en el commit \arch{58b1a9e}. Tras presentar el hito, el MVP se retiró del proyecto (Sección~\ref{sec:arreglos}); su código sigue disponible en el historial de Git.

El MVP tenía cinco pestañas: caso a caso, subir imagen, simulación, modelos y métricas, y estudio de resonancia. Esta última la agregó Juan Manuel Tengonó y fue la base de la aplicación final.

\paragraph{Caso a caso y trazabilidad.} Se elige uno de los 45 casos de prueba (15 por clase) y se ejecutan los cuatro modelos. La Figura~\ref{fig:mvpcaso} muestra un error real: en \arch{caso_04}, con etiqueta Degeneración, los cuatro modelos predicen Protrusión. La Figura~\ref{fig:mvptraza} muestra el recorrido de esa imagen con valores de la ejecución: la ResNet-50 se calcula una sola vez y sus características pasan a los modelos 2, 3 y 4.

\begin{figure}[htbp]
\centering
\includegraphics[width=\linewidth]{mvp_caso.jpg}
\caption{Pestaña Caso a caso del MVP: resultado de los cuatro modelos sobre un caso de prueba.}
\label{fig:mvpcaso}
\end{figure}

\begin{figure}[htbp]
\centering
\includegraphics[width=\linewidth]{mvp_trazabilidad.jpg}
\caption{Trazabilidad de una imagen por los cuatro modelos, con tiempos y valores intermedios.}
\label{fig:mvptraza}
\end{figure}

La pestaña \emph{Simulación} se describe aparte, en la Sección~\ref{sec:simulacion}.

\paragraph{Modelos y métricas.} La pestaña reunía las métricas oficiales, las matrices de confusión, las curvas ROC y el SHA-256 de cada archivo de pesos (Figuras~\ref{fig:mvpmetricas} y~\ref{fig:mvproc}). Esta información ya no está en la aplicación: se conserva en la documentación y en la Sección~\ref{sec:resultados}.

\begin{figure}[htbp]
\centering
\includegraphics[width=\linewidth]{mvp_metricas.jpg}
\caption{Pestaña Modelos y métricas del MVP.}
\label{fig:mvpmetricas}
\end{figure}

\begin{figure}[htbp]
\centering
\includegraphics[width=\linewidth]{mvp_roc.jpg}
\caption{Curvas ROC y barras de exactitud y F1 macro en el MVP.}
\label{fig:mvproc}
\end{figure}

% Capturas tomadas en el commit 58b1a9e, ultimo con el MVP disponible.
% =====================================================================
\section{MVP de demostración}
\label{sec:mvp}

Para el Hito~2 se construyó un MVP local, que se clona y se ejecuta con un solo comando: un backend FastAPI con los cuatro modelos y una interfaz estática en \arch{/mvp}. Su objetivo era demostrar la simulación de resultados y la trazabilidad de una predicción. Las capturas de esta sección se tomaron en el commit \arch{58b1a9e}. Tras presentar el hito, el MVP se retiró del proyecto (Sección~\ref{sec:arreglos}); su código sigue disponible en el historial de Git.

El MVP tenía cinco pestañas: caso a caso, subir imagen, simulación, modelos y métricas, y estudio de resonancia. Esta última la agregó Juan Manuel Tengonó y fue la base de la aplicación final.

\paragraph{Caso a caso y trazabilidad.} Se elige uno de los 45 casos de prueba (15 por clase) y se ejecutan los cuatro modelos. La Figura~\ref{fig:mvpcaso} muestra un error real: en \arch{caso_04}, con etiqueta Degeneración, los cuatro modelos predicen Protrusión. La Figura~\ref{fig:mvptraza} muestra el recorrido de esa imagen con valores de la ejecución: la ResNet-50 se calcula una sola vez y sus características pasan a los modelos 2, 3 y 4.

\begin{figure}[htbp]
\centering
\includegraphics[width=\linewidth]{mvp_caso.jpg}
\caption{Pestaña Caso a caso del MVP: resultado de los cuatro modelos sobre un caso de prueba.}
\label{fig:mvpcaso}
\end{figure}

\begin{figure}[htbp]
\centering
\includegraphics[width=\linewidth]{mvp_trazabilidad.jpg}
\caption{Trazabilidad de una imagen por los cuatro modelos, con tiempos y valores intermedios.}
\label{fig:mvptraza}
\end{figure}

La pestaña \emph{Simulación} se describe aparte, en la Sección~\ref{sec:simulacion}.

\paragraph{Modelos y métricas.} La pestaña reunía las métricas oficiales, las matrices de confusión, las curvas ROC y el SHA-256 de cada archivo de pesos (Figuras~\ref{fig:mvpmetricas} y~\ref{fig:mvproc}). Esta información ya no está en la aplicación: se conserva en la documentación y en la Sección~\ref{sec:resultados}.

\begin{figure}[htbp]
\centering
\includegraphics[width=\linewidth]{mvp_metricas.jpg}
\caption{Pestaña Modelos y métricas del MVP.}
\label{fig:mvpmetricas}
\end{figure}

\begin{figure}[htbp]
\centering
\includegraphics[width=\linewidth]{mvp_roc.jpg}
\caption{Curvas ROC y barras de exactitud y F1 macro en el MVP.}
\label{fig:mvproc}
\end{figure}

% Capturas tomadas en el commit 58b1a9e, ultimo con el MVP disponible.
% =====================================================================
\section{MVP de demostración}
\label{sec:mvp}

Para el Hito~2 se construyó un MVP local, que se clona y se ejecuta con un solo comando: un backend FastAPI con los cuatro modelos y una interfaz estática en \arch{/mvp}. Su objetivo era demostrar la simulación de resultados y la trazabilidad de una predicción. Las capturas de esta sección se tomaron en el commit \arch{58b1a9e}. Tras presentar el hito, el MVP se retiró del proyecto (Sección~\ref{sec:arreglos}); su código sigue disponible en el historial de Git.

El MVP tenía cinco pestañas: caso a caso, subir imagen, simulación, modelos y métricas, y estudio de resonancia. Esta última la agregó Juan Manuel Tengonó y fue la base de la aplicación final.

\paragraph{Caso a caso y trazabilidad.} Se elige uno de los 45 casos de prueba (15 por clase) y se ejecutan los cuatro modelos. La Figura~\ref{fig:mvpcaso} muestra un error real: en \arch{caso_04}, con etiqueta Degeneración, los cuatro modelos predicen Protrusión. La Figura~\ref{fig:mvptraza} muestra el recorrido de esa imagen con valores de la ejecución: la ResNet-50 se calcula una sola vez y sus características pasan a los modelos 2, 3 y 4.

\begin{figure}[htbp]
\centering
\includegraphics[width=\linewidth]{mvp_caso.jpg}
\caption{Pestaña Caso a caso del MVP: resultado de los cuatro modelos sobre un caso de prueba.}
\label{fig:mvpcaso}
\end{figure}

\begin{figure}[htbp]
\centering
\includegraphics[width=\linewidth]{mvp_trazabilidad.jpg}
\caption{Trazabilidad de una imagen por los cuatro modelos, con tiempos y valores intermedios.}
\label{fig:mvptraza}
\end{figure}

La pestaña \emph{Simulación} se describe aparte, en la Sección~\ref{sec:simulacion}.

\paragraph{Modelos y métricas.} La pestaña reunía las métricas oficiales, las matrices de confusión, las curvas ROC y el SHA-256 de cada archivo de pesos (Figuras~\ref{fig:mvpmetricas} y~\ref{fig:mvproc}). Esta información ya no está en la aplicación: se conserva en la documentación y en la Sección~\ref{sec:resultados}.

\begin{figure}[htbp]
\centering
\includegraphics[width=\linewidth]{mvp_metricas.jpg}
\caption{Pestaña Modelos y métricas del MVP.}
\label{fig:mvpmetricas}
\end{figure}

\begin{figure}[htbp]
\centering
\includegraphics[width=\linewidth]{mvp_roc.jpg}
\caption{Curvas ROC y barras de exactitud y F1 macro en el MVP.}
\label{fig:mvproc}
\end{figure}
